\documentclass[11pt]{article}

\usepackage[preprint]{acl}

\usepackage{times}
\usepackage{latexsym}
\usepackage[T1]{fontenc}
\usepackage[utf8]{inputenc}
\usepackage{microtype}
\usepackage{inconsolata}
\usepackage{amsmath}
\usepackage{booktabs}

\newcommand{\tok}[1]{\texttt{\detokenize{#1}}}

\title{Typeling: A 25M-Parameter Language Model and Decoder\\for Offline Typing Prediction}

\author{ScarletKc \\
  Independent researcher \\
  \url{https://github.com/scarletkc/typeling}}

\begin{document}
\maketitle

\begin{abstract}
Keyboards on phones and desktops complete the word being typed and predict the next one, but the models behind
commercial input methods are closed, and open small language models are trained for neither the task nor its
resource budget. We present Typeling, an open system for offline English typing prediction: Typeling-25M, a
25.4M-parameter causal language model obtained by continued pretraining of a small Qwen3-architecture base model on
150M tokens of permissively licensed chat and writing text, and a decoder that turns its token predictions into word
suggestions. The decoder searches the words consistent with the typed letters, repairs typing errors with a
noisy-channel keyboard model scored by the language model, guarantees an exact top-3 candidate list, and mixes in a
5 MB n-gram model and a local user model. On held-out test text, Typeling raises keystroke savings on chat from 34.1\%
to 43.1\% over the strongest n-gram model the same corpus allows, and with every word mistyped it still completes
85.6\% of the mistyped words correctly, against 48.5\% for an n-gram model. Decoder-side typo repair outperformed
fine-tuning the model on synthetic typos, which hurt clean typing. Suggestions take about 25 ms on a laptop CPU. The
model, an 8-bit GGUF file (27.6 MB), the decoder, the benchmark and the training code are released under the Apache
License 2.0.
\end{abstract}

\section{Introduction}

Word prediction is one of the oldest applications of language models: given the text before the cursor, suggest the
rest of the current word and the next word, so that a user types fewer keys
\citep{garay-vitoria-abascal-2006-text}. Mobile keyboards have moved from n-gram models to neural language models
\citep{yu-etal-2018-device,hard-etal-2018-federated,zhang-etal-2024-neural}, and a recent keyboard runs a
fine-tuned 360M-parameter model through llama.cpp on the phone \citep{lexo-2026}. These systems are proprietary. Open
language models, on the other hand, are trained for chat and benchmarks, are hundreds of millions of parameters or
more, and predict subword tokens, not words: used naively for typing, a model asked to continue \texttt{"Let me know
"} produces punctuation, because byte-level tokenizers attach the space to the following word.

We ask how good a typing predictor can be when it is small enough to run anywhere, offline, and built only from
openly licensed parts. Our contributions:

\begin{itemize}
  \item \textbf{Typeling-25M}, a 25.4M-parameter model continued-pretrained for 58 minutes on one laptop GPU from
  SupraLabs/Supra2-Medium-Base \citep{supra2-medium-2026} on chat and writing corpora with permissive licenses, with
  two small changes for typing: a trained text separator and a trailing-space augmentation (\S\ref{sec:model}).
  \item \textbf{A typing decoder} that computes word probabilities from a subword model under the constraint of the
  typed letters, repairs typos with a noisy-channel keyboard model whose candidates share one beam with the typed
  letters, returns an exact top-3 candidate list, and fuses an n-gram model and a user model (\S\ref{sec:decoder}).
  \item \textbf{A benchmark} that simulates typing held-out chat, informal and written text and reports keystroke
  savings for ghost text and for a three-candidate bar, with and without simulated typing errors
  (\S\ref{sec:benchmark}).
  \item \textbf{Findings} from the development of the system (\S\ref{sec:experiments}): an n-gram model saturates
  quickly on a 47M-word corpus; doubling the base model buys only 1--2 points of keystroke savings; continued
  pretraining on 20M tokens closes most of the domain gap; and repairing typos in the decoder beats fine-tuning the
  model on them.
\end{itemize}

Everything is released: the model in 16-bit and 8-bit GGUF form, the Python package with the decoder and a
dependency-free n-gram kernel, the benchmark and the training code.\footnote{\url{https://github.com/scarletkc/typeling},
\url{https://huggingface.co/ScarletKc/Typeling-25M}.}

\section{Related Work}

\paragraph{Word prediction.} Text prediction for augmentative communication and keyboards has a long history
\citep{garay-vitoria-abascal-2006-text}; it is usually evaluated by keystroke savings, whose computation involves
choices (what a keystroke is, how suggestions are shown and accepted) that make numbers hard to compare across
papers \citep{trnka-mccoy-2008-evaluating}. We therefore compare all systems within one simulation and report the
choices explicitly (\S\ref{sec:benchmark}).

\paragraph{Language models in keyboards.} Production keyboards combine a language model with a spatial model of
touches \citep{goodman-etal-2002-language,fowler-etal-2015-effects}. \citet{yu-etal-2018-device} fit a neural word
prediction model into 7.4 MB with 6.5 ms per prediction; \citet{hard-etal-2018-federated} trained a keyboard
language model with federated learning; \citet{zhang-etal-2024-neural} replaced the n-gram language model in the
Gboard decoder's search space with a neural one. \citet{lexo-2026} reports that moving from Qwen2.5-0.5B to 1.5B
raised next-word top-1 accuracy only from 19.0\% to 20.1\% at 2.4 times the decoding cost, the same diminishing
return we find at a much smaller scale. \citet{fowler-etal-2015-effects} measure the benefit of personalizing the
language model, which our decoder supports but our benchmark does not measure.

\paragraph{Noisy-channel correction.} Spelling correction as a noisy channel, choosing the intended word $c$ that
maximizes $P(c)\,P(t \mid c)$ for a typed string $t$, goes back to \citet{kernighan-etal-1990-spelling}, with richer
error models by \citet{brill-moore-2000-improved}. Our typo repair applies the same decomposition to partial words,
with a contextual language model as the prior.

\paragraph{Subword models and character strings.} Language models are distributions over token strings, so
conditioning on a character prefix that ends inside a token biases predictions
\citep{phan-etal-2024-understanding}; \citet{vieira-etal-2024-language} give exact and approximate algorithms to
convert a token-level model into a character-level one. Our decoder solves a narrower version of this problem, word
probabilities under a typed prefix, with a constrained beam search.

\paragraph{Small language models.} Our base model and the references we compare against are small open models:
Supra2-Medium-Base \citep{supra2-medium-2026} in the Qwen3 architecture \citep{qwen3-2025}, Supra-50M-Base, and
Pythia \citep{biderman-etal-2023-pythia}. SmolLM2 \citep{allal-etal-2025-smollm2} starts at 135M parameters, beyond
our size budget.

\section{Task and Benchmark}
\label{sec:benchmark}

\paragraph{Data.} We use five English corpora with permissive licenses (Table~\ref{tab:data}): FineWeb
\citep{penedo-etal-2024-fineweb} for written text; SODA \citep{kim-etal-2023-soda}, everyday dialogue generated by
InstructGPT; the user turns of WildChat-1M \citep{zhao-etal-2024-wildchat} and OpenAssistant oasst2
\citep{kopf-etal-2023-openassistant}, text typed by people into a chat interface; and Ubuntu IRC logs as packaged by
the Common Pile \citep{kandpal-etal-2025-common}, informal internet chat. Emails, URLs and phone numbers are
removed; chat sources keep only turns that look typed by a person. About 1\% of every corpus is held out as a
development split and 1\% as a test split, by content hash. All decisions in this paper were made on the
development splits; the test splits were used once, for the final results (\S\ref{sec:final}). Results are pooled
into \emph{chat} (SODA, WildChat, oasst2), \emph{casual} (IRC) and \emph{writing} (FineWeb).

\begin{table}[t]
\centering
\small
\begin{tabular}{@{}llrr@{}}
\toprule
Corpus & License & Words & Mix \\
\midrule
FineWeb & ODC-By 1.0 & 18.5M & 30\% \\
SODA & CC BY 4.0 & 16.3M & 25\% \\
WildChat-1M (users) & ODC-By 1.0 & 7.5M & 27\% \\
Ubuntu IRC & Public domain & 4.4M & 17\% \\
oasst2 (users) & Apache-2.0 & 0.4M & 1\% \\
\bottomrule
\end{tabular}
\caption{Training corpora (training split) and their share of the continued-pretraining mix.}
\label{tab:data}
\end{table}

\paragraph{Typing simulation.} For every target word $w$ of a held-out text (letters with inner apostrophes), the
simulated user types $i = 0, 1, \dots, |w|-1$ letters, and the system is asked for its top three candidates given
the text before the cursor. Comparisons ignore case. We evaluate two interfaces:
\begin{itemize}
  \item \emph{Ghost text}: the top candidate is shown when its probability is at least a threshold (0.3 for both
  next-word prediction and completion, the defaults of both systems), and the user accepts it with one keystroke as
  soon as it is the target.
  \item \emph{Candidate bar}: the top three candidates are always shown, and the user picks the target with one
  keystroke as soon as it is among them.
\end{itemize}
A word accepted after $i$ letters costs $\min(i+1, |w|)$ keystrokes, an unaccepted word $|w|$. The keystroke savings
rate (KSR) is one minus the keystrokes used over the letters of all target words; KSR@3 is the same for the
candidate bar. We also report the precision of ghost text (accepted over shown suggestions), next-word accuracy
(acc@0: the target is the top candidate before any letter is typed) and top3@0.

\paragraph{Typo simulation.} The simulation is repeated with a typist who types every word of three or more letters
with an error and does not notice it: one error, two with probability 0.2, each one of a dropped apostrophe
(\texttt{dont}), a missing capital, a QWERTY adjacent-key substitution, an omitted letter, an inserted adjacent or
doubled key, or a transposition of neighboring letters. The typist continues through the whole mistyped word;
accepting the intended word corrects the error. Errors depend only on a seed and the texts, so every system sees the
same ones. Besides KSR, we report \emph{Correct}: the share of mistyped words that still end as the intended word.
This is a worst case (76--82\% of all words are mistyped); with a share $r$ of mistyped words, KSR is about $(1-r)$
times the clean KSR plus $r$ times the KSR with typos.

\paragraph{Baselines.} Typeling includes a pure-Python n-gram kernel: an interpolated modified Kneser--Ney trigram
model \citep{kneser-ney-1995-improved,chen-goodman-1999-empirical} over a 60,000-word vocabulary, pruned with the
weighted-difference criterion \citep{seymore-rosenfeld-1996-scalable} to a 5 MB file with one-byte quantized
probabilities, combined with a document cache and a user model. Without pruning, the same corpus yields a 23.4 MB
model (1.66M bigrams, 3.41M trigrams), the strongest n-gram baseline the corpus allows.

\section{Model}
\label{sec:model}

\paragraph{Base model.} We required a recent model under about 50M parameters, under the Apache License, in an
architecture that llama.cpp \citep{llamacpp} supports. Table~\ref{tab:zeroshot} compares candidates zero-shot
through our decoder. Supra-50M-Base (51.8M parameters, Llama architecture) is the most accurate, 1--2 KSR points
ahead of Supra2-Medium-Base (25.4M parameters, Qwen3 architecture: 7 layers, hidden size 512, 8 attention heads,
intermediate size 896, a 16,384-token byte-level BPE vocabulary with tied embeddings, 1,024-token context, trained
on 20B FineWeb-Edu tokens). The gap survives simulated 4-bit quantization of the larger model, which costs it
0.4--1.1 points, consistent with \citet{dettmers-zettlemoyer-2023-case}. We chose Supra2-Medium-Base anyway: it costs
half the compute per token and half the file size for a small loss, and a keyboard model runs next to every other
application. Neither base model had seen chat; both trail the unpruned n-gram model on IRC.

\begin{table}[t]
\centering
\small
\setlength{\tabcolsep}{4pt}
\begin{tabular}{@{}lrrrr@{}}
\toprule
Model & Params & Chat & Casual & Writing \\
\midrule
n-gram & 5 MB & 29.7 & 22.8 & 23.5 \\
n-gram, unpruned & 23 MB & 31.9 & \textbf{24.8} & 25.5 \\
Pythia-14M & 14.1M & 28.5 & 21.8 & 33.6 \\
Pythia-31M & 30.5M & 31.2 & 22.9 & 36.6 \\
Supra2-Medium-Base & 25.4M & 33.5 & 22.8 & 38.6 \\
Supra-50M-Base, Q4$^\dagger$ & 51.8M & 35.0 & 23.4 & 41.0 \\
Supra-50M-Base & 51.8M & \textbf{35.6} & 24.1 & \textbf{42.1} \\
\bottomrule
\end{tabular}
\caption{Zero-shot KSR (\%) of candidate base models through the decoder, development splits, 2,000 words per
source. $^\dagger$Linear layers rounded as GGUF Q4\_0, embeddings to 6 bits.}
\label{tab:zeroshot}
\end{table}

\paragraph{Continued pretraining.} We update all parameters on 150M tokens (about 2.2 epochs of the 67M-token
corpus) with the mix in Table~\ref{tab:data}: packed sequences of 512 tokens, 4,578 steps of 32,768 tokens, AdamW
\citep{loshchilov-hutter-2019-decoupled} with peak learning rate $6\times10^{-4}$, 5\% linear warmup, cosine decay to
zero, weight decay 0.1, bf16 mixed precision. The run took 58 minutes on a laptop RTX 5070 Ti. oasst2 is held at
1\% because at 5\% its 0.5M tokens would be repeated about 15 times. Two changes target typing:
\begin{itemize}
  \item \emph{Separator.} The base model was trained on concatenated documents without separators; its \tok{</s>}
  embedding was untrained, and starting a text with it gave per-word perplexities 3--30\% higher than starting
  with a newline. We separate texts with \tok{</s>} and initialize its embedding with the newline token's.
  \item \emph{Trailing space.} Byte-level BPE attaches the space to the next word (\tok{" know"}), so a text ending
  in a space is rare in training, and the raw model continues \texttt{"Let me know "} with junk such as \tok{"!"}.
  In 1\% of words we re-tokenize \tok{" word"} as \tok{" "} followed by \tok{"word"}. At 3\% the raw model predicts
  sensible words after a space but the development loss rises by about 0.03 on every source and KSR by 0.7 points
  less; at 1\% the loss rises by about 0.01. Our decoder does not need the augmentation (\S\ref{sec:decoder}); it
  helps users of the raw weights.
\end{itemize}

\section{Decoder}
\label{sec:decoder}

The decoder exposes \texttt{topk(text, k)}, the $k$ most likely full words for the text before the cursor, and
\texttt{suggest(text)}, the ghost text.

\paragraph{Words from tokens.} The text is split into the context $c$, the whitespace after it (the lead $\ell$)
and the partially typed word $p$. The context is cut before its trailing whitespace, and the string
$\tau = \ell p$ (e.g.\ \tok{" kn"}) becomes a constraint on the generated tokens. A beam search (width 8) explores
token sequences $y$ in two phases. In the \emph{prefix phase}, only tokens consistent with $\tau$ are allowed: tokens
that are prefixes of the rest of $\tau$ (\tok{" k"}) or that start with it (\tok{" know"}). In the \emph{word phase},
once $\tau$ is covered, the word grows with letter and contraction tokens (\tok{"know"}$+$\tok{"ing"},
\tok{"don"}$+$\tok{"'t"}) until the model ends it. With $C$ the set of word-continuing tokens, the probability that
the word ends after $y$ is $e(y) = 1 - \sum_{t \in C} P(t \mid c, y)$, and
\begin{equation}
P(w \mid c, \tau) = \frac{1}{Z}\sum_{y\,:\,\mathrm{str}(y) = \ell w} P(y \mid c)\, e(y),
\label{eq:word}
\end{equation}
where $Z$, the probability of the constraint, accumulates the mass of all tokens that complete $\tau$ during the
prefix phase. Casings of a word are merged. The context's key/value cache is reused across keystrokes.

\paragraph{Exact top 3.} The search stops when no unexplored path can beat the third-best word found so far. The
first three candidates are therefore exact up to the beam; further candidates are the other words the search
reached.

\paragraph{Typo repair.} For typed letters $p$ (two or more), we enumerate the strings $q$ that one error turns into
$p$: each adjacent-key substitution, each deletion of a key that doubles or neighbors an adjacent letter, each
insertion of a letter or apostrophe (not at the end, where a letter is simply not typed yet), each transposition of
neighboring letters, and capitalizing the first letter. We keep the $q$ that start a word of the n-gram
vocabulary, rank them by the unigram mass of the words they start, and add the best four as extra constraints
$\tau_q = \ell q$ with weight $\omega = 0.003$, the probability of an error relative to none. All constraints share
one beam: each hypothesis starts with $\log \omega_q$ added to its log probability, and hypotheses that satisfy
different constraints merge once their token paths coincide. The probability of a word becomes
\begin{equation}
P(w \mid c, p) \propto \sum_{q \,:\, w \text{ starts with } q} \omega_q\, P(w \mid c),
\end{equation}
normalized by $\sum_q \omega_q P(\tau_q \mid c)$, a noisy channel \citep{kernighan-etal-1990-spelling} with the
language model as prior. The typed letters have weight 1, or 0.1 when they start no vocabulary word: the model has
learned frequent misspellings from chat (\texttt{recieve}) and would otherwise prefer them. Corrections appear in
\texttt{topk()} and the candidate bar; ghost text only extends the typed letters. Because the corrections share the
beam, the number of model calls barely changes.

\paragraph{N-gram fusion and filters.} The final probability is $(1-\lambda)$ times the decoder's plus $\lambda =
0.1$ times the n-gram kernel's (its top 8 words, 0 for others). The short beam misses rare long words
(\texttt{encyc} $\to$ \textit{encyclopedia}: absent without fusion, 0.93 with it), and subword pieces can form
non-words (\textit{don'tknow}); the n-gram vocabulary fixes both. Words the n-gram kernel marks as offensive or as
known misspellings are never suggested.

\paragraph{User model.} The kernel's user model (counts of the words and word pairs the user typed) is mixed into
the joint probabilities before normalization over the words that fit the typed letters:
$P(w \mid c, p) \propto (1-u)\,P_{\mathrm{LM}}(w, p \mid c) + u\,P_{\mathrm{user}}(w \mid w_{-1})$, with
$u = 0.3\,N/(N + 2000)$ growing with the number $N$ of learned words. A name the language model barely knows takes
the lead after a few uses, because the model gives its letters almost no probability, while common words stay with
the language model.

\paragraph{Context.} The decoder uses the last 1,000 characters before the cursor; the start of a cut context moves
in steps of 250 characters so that the cache stays valid in between.

\section{Experiments}
\label{sec:experiments}

\subsection{How far an n-gram model scales}

Removing all pruning (5 MB $\to$ 23 MB) raises development KSR by only 1.8--2.2 points (chat 32.1\% $\to$ 34.3\%,
casual 22.8\% $\to$ 24.6\%, writing 23.1\% $\to$ 25.0\%; 20,000 words per source), and next-word accuracy barely
moves (SODA 28.6\% $\to$ 29.2\%, WildChat 21.3\% $\to$ 22.6\%, IRC 13.8\% $\to$ 14.4\%). A trigram model sees two
words of context; more size only adds rarer word combinations.

\subsection{Continued pretraining}

A 20M-token pilot (5.6 minutes) already raised chat KSR from 33.5\% (zero-shot) to 40.3\% and casual from 22.8\% to
27.7\%, with writing slightly lower (38.6\% $\to$ 37.5\%) as the model moved toward chat. The full 150M-token run
adds 0.9 points on chat (41.2\%) and 1.9 on casual (29.6\%); writing stays flat (37.2\%), and development loss per
token falls from 2.85 to 1.71 on SODA, 4.49 to 3.52 on WildChat and 4.72 to 3.82 on IRC, while FineWeb rises from
3.75 to 3.79. Returns diminish on this corpus: more data, not more epochs, is the likelier next gain. A higher
learning rate ($1.2\times10^{-3}$) adapted slightly faster to chat (development loss 0.01--0.02 lower) but forgot
more writing (0.10 higher), and $3\times10^{-4}$ the reverse.

\subsection{Final results}
\label{sec:final}

Table~\ref{tab:final} reports the test splits, run once with the released 16-bit weights and the decoder defaults.
Typeling-25M saves 9.0 points more keystrokes than the unpruned n-gram model on chat, 5.3 on casual and 8.5 on
writing, at higher precision; the candidate bar adds 15--18 points over ghost text for every system. Next-word
accuracy (Table~\ref{tab:nextword}) is 5--13 points higher on every source.

\begin{table*}[t]
\centering
\small
\begin{tabular}{@{}lccccccc@{}}
\toprule
& \multicolumn{2}{c}{Chat} & \multicolumn{2}{c}{Casual} & \multicolumn{2}{c}{Writing} & Precision \\
\cmidrule(lr){2-3}\cmidrule(lr){4-5}\cmidrule(lr){6-7}
System & KSR & KSR@3 & KSR & KSR@3 & KSR & KSR@3 & (chat) \\
\midrule
n-gram, 5 MB & 31.6 & 48.8 & 22.7 & 38.9 & 22.3 & 40.6 & 51.0 \\
n-gram, unpruned, 23 MB & 34.1 & 50.6 & 24.4 & 40.4 & 23.9 & 41.5 & 51.2 \\
Typeling-25M & \textbf{43.1} & \textbf{58.0} & \textbf{29.8} & \textbf{46.8} & \textbf{32.4} & \textbf{49.9} &
\textbf{54.7} \\
\bottomrule
\end{tabular}
\caption{Test splits, 2,000 words per source: keystroke savings (\%) with ghost text (KSR) and a three-candidate bar
(KSR@3), and ghost-text precision on chat.}
\label{tab:final}
\end{table*}

\begin{table}[t]
\centering
\small
\begin{tabular}{@{}lcc@{}}
\toprule
Source & n-gram, 23 MB & Typeling-25M \\
\midrule
SODA & 30.0 / 46.9 & \textbf{43.4 / 62.5} \\
WildChat & 22.6 / 35.5 & \textbf{31.3 / 45.9} \\
oasst2 & 18.9 / 30.7 & \textbf{28.1 / 44.5} \\
Ubuntu IRC & 13.9 / 24.0 & \textbf{19.0 / 32.4} \\
FineWeb & 15.0 / 27.1 & \textbf{24.8 / 37.7} \\
\bottomrule
\end{tabular}
\caption{Next word before any letter is typed, test splits: acc@0 / top3@0 (\%).}
\label{tab:nextword}
\end{table}

\subsection{Typos: repair in the decoder or in the model?}

Typing every word with an error costs the decoder without repair 8.4--8.8 KSR points on the development splits
(Table~\ref{tab:typos}). Repair recovers 36--48\% of that loss at a clean cost of 0.1--0.5 points and raises the
share of mistyped words that end correct from 48--62\% to 72--84\%. The error weight matters little between 0.003
and 0.03; higher weights cost more on clean input.

We also fine-tuned the model on typos: 40M tokens from the continued-pretraining checkpoint, half the usual corpus
mix and half 600,000 examples of the form \tok{context <|typed|> letters <|word|> word </s>} with the loss on the
intended word only (47\% of the typed letters carry errors). The typo loss fell from 5.78 to 1.14 per token and the
model corrects well in isolation (the intended word on top after the whole mistyped word for 84\% of SODA words,
against 7\% before), but clean development loss stayed 0.07--0.12 above the checkpoint, and in the typing simulation
it is worse than decoder repair everywhere: 6--10 points lower KSR on clean text and 6--11 on mistyped text. We
therefore release the model without typo fine-tuning.

On the test splits, with every word mistyped, Typeling's chat KSR falls from 44.5\% to 40.1\% (85.6\% of mistyped
words end correct) and KSR@3 from 59.4\% to 55.7\% (94.5\%), against 32.1\% $\to$ 22.7\% (48.5\%) and 49.6\% $\to$
40.5\% (71.0\%) for the 5 MB n-gram model; casual and writing behave alike (Appendix~\ref{app:typos}).

\begin{table}[t]
\centering
\footnotesize
\setlength{\tabcolsep}{3pt}
\begin{tabular}{@{}lccc@{}}
\toprule
Decoder & Chat & Casual & Writing \\
\midrule
n-gram, 5 MB & 29.4 $\to$ 20.4 & 21.6 $\to$ 14.3 & 22.7 $\to$ 14.7 \\
& (48.4) & (39.2) & (39.2) \\
No repair & 41.3 $\to$ 32.5 & 28.8 $\to$ 20.2 & 34.7 $\to$ 26.3 \\
& (61.7) & (47.9) & (54.3) \\
Repair, $\omega{=}0.003$ & 41.1 $\to$ \textbf{36.5} & 28.3 $\to$ \textbf{22.8} & 34.6 $\to$ \textbf{29.6} \\
& (84.3) & (72.1) & (82.0) \\
Repair, $\omega{=}0.03$ & 40.7 $\to$ 36.7 & 27.6 $\to$ 22.7 & 33.6 $\to$ 29.6 \\
& (85.5) & (74.7) & (82.8) \\
Typo fine-tuned & 34.0 $\to$ 28.8 & 22.6 $\to$ 16.7 & 24.5 $\to$ 19.1 \\
& (73.9) & (62.5) & (61.4) \\
\bottomrule
\end{tabular}
\caption{Development splits, 1,000 words per source: KSR (\%) clean $\to$ with every word mistyped, and in
parentheses the share of mistyped words that end correct. Without n-gram fusion.}
\label{tab:typos}
\end{table}

\subsection{Decoder ablations}

On the development splits (1,000 words per source), n-gram fusion at $\lambda = 0.1$ changes KSR@3 by $+0.2$ to
$+0.3$ points and ghost-text precision by $+1$ to $+2$, and lowers ghost-text KSR by up to 0.4; $\lambda = 0.3$ lowers
chat KSR by 0.8. Down-weighting typed letters that start no vocabulary word changes nothing on clean text and adds
0.1--0.3 points with typos. These differences are within sampling noise (\S\ref{sec:limitations}); we keep both
for the qualitative errors they fix (\S\ref{sec:decoder}). A context of 1,000 instead of 2,000 characters costs 0.1
KSR points on writing and nothing on chat, whose texts are shorter, and lowers the 95th and 99th percentile CPU
time on long texts from 76 and 110 ms to 63 and 82 ms.
A beam of 6 is 0.1--0.6 points worse than 8 and not faster.

\subsection{Export and quantization}

The 16-bit safetensors (50.8 MB) give the same development perplexity as float32. The Q8\_0 GGUF file (27.6 MB),
compared with an F16 conversion on 380k characters of development text, raises perplexity by 0.08\% (30.025 vs
30.000), with a mean KL divergence of 0.0009 and the same top token 97.9\% of the time. We release no 4-bit file: at
25M parameters it would save about 10 MB, and small models lose more to 4 bits.

\subsection{Latency and footprint}

Table~\ref{tab:latency} reports the time per suggestion of the full decoder in PyTorch on a laptop CPU (AMD Ryzen 9
9955HX, 4 threads, float32). Chat and casual text meet our target of 50 ms at the 95th percentile and 100 ms at the
99th, the gap between keystrokes of a fast typist; long written text meets only the second. Most of the time goes
into copying the context's key/value cache for every hypothesis at every beam step; a llama.cpp implementation, whose
sequences can share cache cells, would remove it. On the same 4 threads, llama.cpp processes 11,900 tokens per second
of prompt and 2,800 single-token steps per second with the Q8\_0 file. The PyTorch path occupies 549 MB of memory
(the n-gram kernel 54 MB) and needs 4.6 s to import PyTorch and 2.5 s to load the model, acceptable for a library
but not for a resident input method.

\begin{table}[t]
\centering
\small
\begin{tabular}{@{}lccc@{}}
\toprule
Source & Mean & p95 & p99 \\
\midrule
SODA & 24.8 & 39.1 & 51.1 \\
WildChat & 23.7 & 39.9 & 52.0 \\
oasst2 & 23.5 & 44.7 & 56.0 \\
Ubuntu IRC & 23.4 & 47.7 & 59.7 \\
FineWeb & 32.8 & 63.0 & 82.2 \\
\bottomrule
\end{tabular}
\caption{Milliseconds per suggestion, PyTorch decoder with typo repair and n-gram fusion on a laptop CPU (4
threads), 300 words per development source. The n-gram kernel alone takes 41--115~$\mu$s.}
\label{tab:latency}
\end{table}

\section{Conclusion}

A 25M-parameter model, adapted to chat in under an hour on a laptop and decoded with care for words, typos and the
typed prefix, saves 9 points more keystrokes on chat than the strongest n-gram model its corpus supports and repairs
most typing errors. The largest remaining costs are engineering ones: the memory and startup time of PyTorch and the
latency of long contexts, which a llama.cpp port of the decoder should address.

\section*{Limitations}
\label{sec:limitations}

\begin{itemize}
  \item \textbf{No confidence intervals.} Each cell uses 1,000--2,000 words per source (300 for latency); from
  repeated runs, differences of about one point are within sampling noise. The main gaps (5--9 points) are well
  beyond it; the decoder ablations (0.1--0.4 points) are not.
  \item \textbf{Quantized file not run through the typing benchmark.} The Q8\_0 GGUF file was compared with the
  16-bit weights by perplexity and KL divergence only, because the decoder runs on PyTorch.
  \item \textbf{No comparison with larger general models.} We did not evaluate models of hundreds of millions of
  parameters (e.g.\ SmolLM2-360M, Qwen2.5-0.5B) through the same decoder; the only evidence on scale beyond 50M
  parameters is the report of \citet{lexo-2026}.
  \item \textbf{Simulated users.} Typing and typos are simulated; the typo model reflects common keyboard errors,
  not recorded human typing, and no user study was run. The benchmark cannot measure personalization, because its
  texts come from many different people.
  \item \textbf{Data.} SODA was generated by InstructGPT, whose terms restrict using outputs to build competing
  models. The base model was trained on FineWeb-Edu, filtered from FineWeb, so writing results may be optimistic.
  English only.
  \item \textbf{One machine.} Latency was measured on one laptop CPU.
\end{itemize}

\section*{Ethical Considerations}

Typeling runs offline: what the user types, and the user model learned from it, stay on the machine. All training
data carries permissive licenses and is attributed in the release. The model can suggest offensive words or
reproduce biases of web and chat text; the decoder does not suggest words from a blocklist, but other harmful
completions are possible.

\bibliography{references}

\appendix

\section{Typos on the Test Splits}
\label{app:typos}

Test splits, 1,000 words per source: keystroke savings (\%) clean $\to$ with every word of three or more letters
mistyped, and in parentheses the share of mistyped words that end correct.

\begin{center}
\footnotesize
\setlength{\tabcolsep}{3pt}
\begin{tabular}{@{}llcc@{}}
\toprule
& & n-gram, 5 MB & Typeling-25M \\
\midrule
Chat & KSR & 32.1 $\to$ 22.7 (48.5) & 44.5 $\to$ 40.1 (85.6) \\
& KSR@3 & 49.6 $\to$ 40.5 (71.0) & 59.4 $\to$ 55.7 (94.5) \\
Casual & KSR & 22.5 $\to$ 14.8 (36.9) & 28.6 $\to$ 23.1 (73.1) \\
& KSR@3 & 39.1 $\to$ 29.3 (57.8) & 46.5 $\to$ 41.3 (88.4) \\
Writing & KSR & 24.3 $\to$ 15.7 (40.9) & 33.7 $\to$ 27.9 (78.5) \\
& KSR@3 & 42.3 $\to$ 32.0 (62.9) & 50.8 $\to$ 45.5 (90.5) \\
\bottomrule
\end{tabular}
\end{center}

\section{Examples}

Suggestions of the released model with the default decoder (top candidates with probabilities):
\begin{itemize}
  \item \texttt{"Let me know "}: \textit{if} 0.73, \textit{how} 0.12, \textit{what} 0.06.
  \item \texttt{"See you tomm"}: \textit{tomorrow} 0.90, \textit{tommy} 0.07.
  \item \texttt{"Thank you so mich"}: \textit{much} 0.90, \textit{michael} 0.04.
  \item \texttt{"What do you thnik"}: \textit{think} 0.89.
  \item \texttt{"I looked it up in the encyc"}: \textit{encyclopedia} 0.93.
  \item \texttt{"I will recie"}: \textit{receive} first, but at 0.01: the model's mass goes to the misspelling
  \textit{recieve}, which the decoder hides, so no ghost text is shown.
\end{itemize}

\section{Reproduction}

The repository contains the corpus builder (\texttt{training/build.py}), the tokenizer and training scripts
(\texttt{lm/data.py}, \texttt{lm/train.py}), the benchmark (\texttt{bench/eval.py}, \texttt{lm/eval.py}), the
export to safetensors and GGUF (\texttt{lm/export.py}), and \texttt{docs/model-plan.md}, which records every run
in this paper with its command and results; \texttt{bench/results.jsonl} stores the raw benchmark output.

\end{document}
